\begin{chineseabstract}
水空直接无线光通信使水下节点能够不经海面中继而与空中平台直接交换数据。上行光束从海水进入空气时，裸水面在 $48.75^\circ$ 处发生全反射，亚临界斜入射的光束被折射到接近掠射的方向；风浪使海面起伏、海面平台随波运动，收发两端的位置抖动、指向误差和湍流衰落进一步降低了链路可靠性。本文在海面上漂浮一块贴水安装的透射式液晶智能超表面（RIS），在光束进入空气之前对其施加逐单元可编程相位，围绕坐标与器件、平台运动与相位设计、入射与接收端复光场以及动态海面下的链路性能建立系统模型，并评估 RIS 对链路性能的增强作用。

首先，建立了全文统一的坐标体系，在均匀水体、高斯光源和吸收型 RIS 边界下给出辐射输运方程的碰撞级数解，把分层介质中的 Fermat 路径、面积 Jacobian 与 Fresnel 透射组合为复光场传递关系，以 Neumann 方向谱生成随机海面；设计了 $1\,\mathrm m\times1\,\mathrm m$、$40\times40$ 单元的玻璃夹层液晶 RIS，由 1\,$\mu$m 电极周期得到 $15.426^\circ$ 的单轴偏转角上限与 $0.021^\circ$ 的量化步长，论证了贴水安装使全反射的起始入射角推迟到 $72.15^\circ$。

其次，完成了浮框配平，建立了耦合薄板辐射/绕射载荷、侧壁压力解析偏航力矩和被动辐射记忆的六自由度刚体运动模型，得到风速 3--12\,m/s 下 RIS 的姿态与轨迹；在给定位姿下，以全部单元的公共可达方向确定接收机位置，以入射能量的白化径向分位数分配平顶或鲁棒目标落点，经分层 Fermat 路径反解、限幅量化和常相位校准得到可加载相位，并提出预测调相与鲁棒光斑结合的调相方式。

再次，将 RIS 下表面入射光强分解为弹道分量与完整 Henyey--Greenstein 散射分量，在 18 个冻结快照上与独立蒙特卡洛结果的差异处于后者统计噪声水平；对任意可加载相位，以逐单元衍射积分、Chirp-$z$ 变换、功率守恒搬运和像素内相干平均建立 RIS 调制后接收端的复光场模型，在 10--50\,m 和 3--12\,m/s 的 228 个相位上与独立蒙特卡洛结果的接收功率差异在 10\% 以内。

最后，在冻结的动态粗糙海面下，以蒙特卡洛计算无 RIS 直达链路，与 Steering-only、Flat top、Robust ABER/OP 和 Max mean power 四个 RIS 方案共用同一个抖动接收机进行比较。结果表明：接收距离 30\,m 及以上时，直达链路在 90\,dBm 以内达不到 $10^{-3}$ 的平均误码率，最优 RIS 方案只需约 42\,dBm，且所需功率几乎不随接收距离和风速变化；Robust ABER/OP 与 Flat top 是明显优于其他方案的两个方法，在 30\,m 处比 Steering-only 低约 6.5\,dB，比最大平均功率方案低约 14\,dB，鲁棒光斑在强海洋湍流下的中断概率上进一步领先。本文的模型与结果为水空直接无线光通信系统中 RIS 的器件选型、相位设计与链路评估提供了依据。

\chinesekeyword{水空无线光通信，智能超表面，液晶相位调制，六自由度运动，复光场，鲁棒光斑}
\end{chineseabstract}
